\section{构网型电池储能模型}
\label{sec:battery}

电池储能承担小时尺度的能量搬移，并为构网运行预留一定功率和能量裕度。设$P_t^{\mathrm{ch}}$、$P_t^{\mathrm{dis}}$分别为母线侧充电和放电功率，$E_t^{\mathrm B}$为时段$t$结束时的可用能量，则
\begin{equation}
E_t^{\mathrm B}
=\rho_{\mathrm B}E_{t-1}^{\mathrm B}
+\eta_{\mathrm{ch}}P_t^{\mathrm{ch}}\Delta t
-\frac{P_t^{\mathrm{dis}}\Delta t}{\eta_{\mathrm{dis}}}.
\label{eq:bess-state}
\end{equation}
式中$\rho_{\mathrm B}$为能量保持率，$\eta_{\mathrm{ch}}$和$\eta_{\mathrm{dis}}$分别为充、放电效率。充电时只有$\eta_{\mathrm{ch}}P_t^{\mathrm{ch}}\Delta t$进入电池；向母线输出$P_t^{\mathrm{dis}}\Delta t$时，电池内部需释放$P_t^{\mathrm{dis}}\Delta t/\eta_{\mathrm{dis}}$。

SOC和功率边界写为
\begin{equation}
\begin{gathered}
E^{\mathrm r}SOC_{\min}\le E_t^{\mathrm B}\le E^{\mathrm r}SOC_{\max},\\
0\le P_t^{\mathrm{ch}}\le A_t^{\mathrm B}P^{\mathrm{ch,max}},\qquad
0\le P_t^{\mathrm{dis}}\le A_t^{\mathrm B}P^{\mathrm{dis,max}}.
\end{gathered}
\label{eq:bess-soc}
\end{equation}
其中$E^{\mathrm r}$为额定能量，$A_t^{\mathrm B}$为设备可用系数。为避免同一时段同时充、放电，引入$u_t^{\mathrm{ch}},u_t^{\mathrm{dis}}\in\{0,1\}$：
\begin{equation}
P_t^{\mathrm{ch}}\le u_t^{\mathrm{ch}}P^{\mathrm{ch,max}},\quad
P_t^{\mathrm{dis}}\le u_t^{\mathrm{dis}}P^{\mathrm{dis,max}},\quad
u_t^{\mathrm{ch}}+u_t^{\mathrm{dis}}\le1.
\label{eq:bess-exclusive}
\end{equation}

构网支撑需要在常规调度之外保留调节空间。令$g_t\in\{0,1\}$表示时段$t$是否承担构网支撑，$E^{\mathrm{GFM,up}}$和$E^{\mathrm{GFM,down}}$为向上、向下调节所需的能量裕度，则
\begin{equation}
E^{\mathrm r}SOC_{\min}+g_tE^{\mathrm{GFM,up}}
\le E_t^{\mathrm B}
\le E^{\mathrm r}SOC_{\max}-g_tE^{\mathrm{GFM,down}}.
\label{eq:gfm-energy-reserve}
\end{equation}
对应的功率裕度为
\begin{equation}
\begin{aligned}
(P_t^{\mathrm{dis}}-P_t^{\mathrm{ch}})+g_tP^{\mathrm{GFM,up}}
&\le P^{\mathrm{dis,max}},\\
(P_t^{\mathrm{ch}}-P_t^{\mathrm{dis}})+g_tP^{\mathrm{GFM,down}}
&\le P^{\mathrm{ch,max}}.
\end{aligned}
\label{eq:gfm-up-reserve}
\end{equation}
$P^{\mathrm{GFM,up}}$和$P^{\mathrm{GFM,down}}$分别为预留的向上、向下有功调节能力。上述约束限定小时级调度工作点；频率、电压、故障穿越和黑启动等快速动态性能在相应工作点上另行校核。
